\subsection{ADJOINT EQUATION}

Assuming always that the space E is of finite dimension n over C, let \(E^{*}\) be its dual (A, II, p.~40), which is a space of dimension n over C (Alg., II, p.~299, th. 4); the canonical bilinear form \(\langle x, x^{*}\rangle\) defined on \(E \times E^{*}\) (Alg., II, p.~234) is continuous on this product (being a polynomial in the components of \(x \in E\) and \(x^{*} \in E^{*}\) ).

Given a homogeneous linear equation (4) (IV, p.~183), where \(t \mapsto A(t)\) is a regulated map of \(\mathbf{J}\) into \(\mathcal{L}(\mathrm{E})\) , let us see if there exists a map \(t \mapsto \mathbf{v}(t)\) of \(\mathbf{J}\) into \(\mathrm{E}^*\) , a primitive of a regulated function on \(\mathbf{J}\) , and such that the scalar function \(t \mapsto \langle \mathbf{u}(t), \mathbf{v}(t) \rangle\) is constant on \(\mathbf{J}\) when \(\mathbf{u}\) is an arbitrary solution of (4); it comes to the same to write that the derivative of this function should be zero at every point where \(\mathbf{u}\) and \(\mathbf{v}\) are differentiable, that is, one must have

\[
\left\langle \frac {d {\bf u}}{d t}, {\bf v} (t) \right\rangle + \left\langle {\bf u} (t), \frac {d {\bf v}}{d t} \right\rangle = 0
\]

at such points.

Now, by (4), \(\left\langle\frac{d\mathbf{u}}{dt},\mathbf{v}(t)\right\rangle=\langle A(t).\mathbf{u}(t),\mathbf{v}(t)\rangle=-\langle\mathbf{u}(t),B(t).\mathbf{v}(t)\rangle\) where \(-B(t)\) is the transpose of \(A(t)\) (Alg., II, p.~234). The relation that v must satisfy can thus be written

\[
\left\langle \mathbf {u} (t), \frac {d \mathbf {v}}{d t} - B (t). \mathbf {v} (t) \right\rangle = 0
\]

at all points where \(A(t)\) is continuous and \(\mathbf{v}(t)\) is differentiable. Now for such a point \(t\) and an arbitrary point \(\mathbf{x}_0 \in E\) there exists, by th. 1 of IV, p.~179, a solution \(\mathbf{u}\) of (4) such that \(\mathbf{u}(t) = \mathbf{x}_0\) ; one thus must have \(\left\langle \mathbf{x}_0, \frac{d\mathbf{v}}{dt} - B(t).\mathbf{v}(t) \right\rangle = 0\) for all \(\mathbf{x}_0 \in E\) , which entails \(\frac{d\mathbf{v}}{dt} - B(t).\mathbf{v}(t) = 0\) . Consequently:

PROPOSITION 6. The map \(t \mapsto \mathbf{v}(t)\) of \(\mathbf{J}\) into \(\mathrm{E}^*\) , a primitive of a regulated function on \(\mathbf{J}\) , is such that \(\langle \mathbf{u}(t), \mathbf{v}(t) \rangle\) is constant on \(\mathbf{J}\) for every solution \(\mathbf{u}\) of the equation (4) of IV, p.~183 if and only if \(\mathbf{v}\) is a solution of the homogeneous linear equation

\[
{\frac {d \mathbf {x}}{d t}} = B (t). \mathbf {x}\tag{16}
\]

where \(-B(t)\) is the transpose of \(A(t)\) .

Equation (16) is called the adjoint of (4); clearly (4) is the adjoint of (16). The elements of the matrix \(B(t)\) being regulated functions of t on J, the results obtained above on linear equations apply to (16). In particular, the integral of (16) taking the value \(x_{0}^{*}\) at the point \(t_{0}\) can be written as \(H(t, t_{0})\) . \(x_{0}^{*}\) , where \(H(t, t_{0})\) is a bijective linear map of \(E^{*}\) onto itself, identical to the integral of the equation

\[
\frac {d V}{d t} = B (t) V\tag{17}
\]

which takes the value \(I\) at the point \(t_0\) . As a result one has (with the notation of IV, p.~179)

\[
\left\langle C (t, t _ {0}). \mathbf {x} _ {0}, H (t, t _ {0}). \mathbf {x} _ {0} ^ {*} \right\rangle = \left\langle \mathbf {x} _ {0}, \mathbf {x} _ {0} ^ {*} \right\rangle
\]

for any \(x_{0} \in E\) and \(x_{0}^{*} \in E^{*}\) , which shows that

\[
H (t, t _ {0}) = \check {C} (t, t _ {0})\tag{18}
\]

(the contragradient of \(C(t, t_{0})\) ). In particular, if one knows a fundamental system of integrals of the adjoint equation (16) the matrix \(H(t, t_{0})\) is determined, so is \(C(t, t_{0})\) also, and as a result, so are all the integrals of equation (4).

Remark. Let E and F be two complete normed spaces over R (or over C), and \((\mathbf{x}, \mathbf{y}) \mapsto \langle \mathbf{x}, \mathbf{y} \rangle\) a continuous bilinear form on \(E \times F\) , such that the relation `` \(\langle x, y \rangle = 0\) for all \(y \in F\) '' (resp. \(\langle x, y \rangle = 0\) for all \(x \in E\) '') implies x = 0 (resp. y = 0). Suppose further that for every \(t \in J\) there is a continuous linear map \(B(t)\) of F into itself, such that \(\langle A(t).x, y \rangle + \langle x, B(t).y \rangle = 0\) for every \((\mathbf{x}, \mathbf{y}) \in E \times F\) . In these circumstances one sees as before that for a map \(t \mapsto v(t)\) of J into F, a primitive of a regulated function, to be such that \(\langle u(t), v(t) \rangle\) is constant for every integral u of (4), it is necessary and sufficient that v should be an integral of (16), which again we call the adjoint of (4).

